\section{Limite locale et classification modulaire du système infini}

Dans une UHF infinie, il n’existe généralement pas de matrice densité globale à trace à l’intérieur de l’algèbre. L’objet correct est le réseau
\begin{equation}
\label{rm:eq:modular-local-net}
\HHM^{(m)}=-\log\rho_m
\end{equation}
avec des restrictions d’états compatibles, ou le générateur modulaire dans la représentation GNS. L’espérance traciale \(\mathrm{id}\otimes\tau_9\) ne préserve pas automatiquement un état de Gibbs non tracial et ne se substitue pas à la compatibilité des états.

\begin{theorem}[Type modulaire sous persistance fidèle des deux canaux]
\label{rm:thm:modular-typeIII}
Supposons que :
\begin{enumerate}[label=(\roman*)]
\item chaque coupe est un produit d’états qutrit fidèles de poids \(w_{90}\) et \(w_{120}\) ;
\item les inclusions préservent l’état et l’étiquette de canal ;
\item une infinité de facteurs de chaque type apparaît, les deux types ayant une densité asymptotique positive ;
\item aucun quotient spectral supplémentaire n’est introduit et la représentation produit est factorielle.
\end{enumerate}
Alors le groupe additif des logarithmes des quotients de valeurs propres est
\begin{equation}
\label{rm:eq:modular-ratio-group}
90\Dmod\Z+120\Dmod\Z=30\Dmod\Z.
\end{equation}
La fermeture GNS ITPFI est le facteur hyperfini
\begin{equation}
\label{rm:eq:modular-typeIII-lambda}
\boxed{
III_{\lambda_\partial},
\qquad
\lambda_\partial=\ee^{-30\Dmod}
=0.9966696464689573786108299769\ldots .}
\end{equation}
\end{theorem}

\begin{proof}
Dans chaque lien, les quotients de valeurs propres sont des puissances de \(\ee^{-w_m}\). La persistance infinie des deux types engendre le sous-groupe \(w_{90}\Z+w_{120}\Z\). Comme \(\gcd(90,120)=30\), on obtient \cref{rm:eq:modular-ratio-group}. Le critère de rapport asymptotique d’Araki--Woods pour le produit fidèle donne le facteur ITPFI hyperfini de paramètre \(\ee^{-30\Dmod}\) \cite{rm:ArakiWoods1968,rm:Takesaki}.
\end{proof}

Le théorème déclare ses hypothèses effectives. Si l’un des canaux disparaît, si le produit n’est pas compatible ou si le groupe des rapports change, le type doit être recalculé.
